\pdfoutput=1
\documentclass[11pt]{article}
\usepackage{mathfin-paper}

\title{A Machine-Checked It\^o Calculus\\ for Brownian Motion}
\author{Raphael Coelho\thanks{Independent researcher. ORCID
  \href{https://orcid.org/0009-0001-6601-1023}{0009-0001-6601-1023}.
  Correspondence: \href{mailto:raphaelrrcoelho@gmail.com}{\texttt{raphaelrrcoelho@gmail.com}}.
  Artifact: \url{https://github.com/formal-applied-math/formal-mathfin}, release \texttt{v1.5.0}.}}
\date{October 2026}

\hypersetup{pdftitle={A Machine-Checked Ito Calculus for Brownian Motion},
  pdfauthor={Raphael Coelho}}

\newcommand{\HT}{H_T}
\newcommand{\bulletB}{\mathbin{\bullet}}

\begin{document}
\maketitle

\begin{abstract}
\noindent
We develop the It\^o calculus of Brownian motion, machine-checked in Lean~4 over Mathlib and the
\texttt{BrownianMotion} package. On a bounded interval $[0,T]$ the It\^o integral is built as a
Hilbert-space isometry, from a predictable-rectangle $\pi$-system through the density of simple adapted
processes. One identity, that the integral at time $t$ is the conditional-expectation projection of its
terminal value onto $\F_t$, makes it an $L^2$ martingale, continuous in $L^2$, with its time-indexed
isometry; a maximal inequality then gives it a modification with continuous paths, a local martingale on the
half-line. Two theorems pin the
operator down. Its range together with the constants is all of $L^2(\F^B_T)$, which is martingale
representation, proved by totality of Dol\'eans exponentials with no Malliavin calculus; and it
satisfies the \texttt{BrownianMotion} package's axiomatic characterisation of a stochastic integral.
It\^o's formula for $f(t,B_t)$ is proved in $L^2$ for $f\in C^3$ whose derivatives grow at most
exponentially, with the stochastic term the It\^o integral of the named integrand
$\partial_xf(s,B_s)$; for every $f\in C^3$, with no growth condition, as a continuous local martingale,
by localization at exit times; and at a fixed time for It\^o processes with bounded adapted continuous
coefficients and $f\in C^2$ with bounded derivatives. To our knowledge these are the first machine-checked constructions of the It\^o
integral and proofs of It\^o's formula in any proof assistant. The boundary is explicit: the
integrator is Brownian motion or an It\^o integral against it, integrands are square-integrable, and
quadratic variation is established in mean, not pathwise. The development is about 27,000 lines of Lean
in 81 modules, \texttt{sorry}-free, with the axioms of each headline result pinned to Lean's three
standard axioms by a build-enforced gate.
\end{abstract}

%======================================================================
\section{Introduction}

Stochastic calculus sits oddly in the formalization landscape. Its prerequisites (measure-theoretic
probability, conditional expectation, $L^p$ spaces, Gaussian measures) are well represented in the major
proof-assistant libraries. Its consumers (mathematical finance, stochastic control, stochastic partial
differential equations) are fields with strong correctness cultures that cite formal methods as an
aspiration. Yet the calculus itself long remained unformalized. Martingales in discrete and continuous
time were formalized in Isabelle/HOL by Keskin~\cite{Keskin2023}; the central limit theorem was
machine-checked a decade ago~\cite{AvigadHolzlSerafin2017}; and in 2025 Degenne, Ledvinka, Marion and
Pfaffelhuber completed the first formal construction of Brownian motion, in
Lean~4~\cite{DegenneEtAl2025,BrownianMotionRepo}. Their package is the foundation this work builds on,
and it is now growing a general theory of stochastic integration: an axiomatic characterisation of the
stochastic integral, local martingales, and the Doob--Meyer decomposition. The It\^o integral of
Brownian motion, It\^o's formula and the martingale theory of the integral are the layer this paper
supplies.

The development is the stochastic-calculus core of a larger library of formalized mathematical
finance~\cite{CoelhoLibrary}, in which it carries the continuous-time results: Girsanov's theorem,
martingale representation and market completeness, stochastic differential equations, and the
Black--Scholes pricing that rests on them; the library's no-arbitrage theory is treated
separately~\cite{CoelhoFTAP}. Where the library overview surveys breadth, this paper is for readers who
want to know exactly what was proved about the calculus, under which hypotheses, and by which proof
architecture.

\paragraph{Setting}
All results are stated over a probability space $(\Omega,\F,\mu)$ carrying a process
$B\colon\R_{\ge0}\to\Omega\to\R$ whose finite-dimensional laws are those of Brownian motion (the
pre-Brownian interface \lean{IsPreBrownianReal}, developed in the \texttt{BrownianMotion} package and now
part of Mathlib), with measurable coordinates and, where stated, continuous paths. $(\F_t)$ is the natural filtration of $B$
(\lean{natFiltration}).

\paragraph{Contributions}
\begin{enumerate}[label=(\arabic*)]
\item \emph{The It\^o integral as an isometry} (\cref{sec:integral}). The integrand space $\HT$ is the
  $L^2$ space of the product measure on $(0,T]\times\Omega$ trimmed to the predictable $\sigma$-algebra;
  simple adapted processes are dense in it, and the It\^o integral $\ito_T$ (\lean{itoIntegralCLM_T}) is
  the unique continuous linear extension of the elementary integral, an isometry.
\item \emph{The integral as a process} (\cref{sec:process}). The load-bearing identity
  $(\varphi\bulletB B)_t=\E[\ito_T(\varphi)\mid\F_t]$ in $L^2$ makes adaptedness, the martingale property,
  the contraction bound, the terminal and time-indexed It\^o isometries and $L^2$-continuity in $t$
  corollaries.
\item \emph{It\^o's formula in $L^2$} (\cref{sec:ito}). For $f(t,x)$ whose partial derivatives
  $f_t,f_x,f_{xx},f_{tt},f_{tx},f_{xxx}$ are bounded, and then for those growing at most like
  $e^{\lambda\abs{x}}$,
  \[
    f(T,B_T)-f(0,B_0)=\ito_T\bigl(f_x(\cdot,B)\bigr)+\int_0^T\bigl(f_t+\tfrac12f_{xx}\bigr)(s,B_s)\,\dd s
    \quad\mu\text{-a.s.},
  \]
  with the integrand of $\ito_T$ named as the class of $(s,\omega)\mapsto f_x(s,B_s(\omega))$, by a
  discrete-to-continuous argument through weighted quadratic variation with explicit $L^2$ remainder
  bounds and, for the growth case, a cutoff limit carried by the isometry.
\item \emph{The pathwise integral} (\cref{sec:pathwise}). The integral process has an almost surely
  continuous modification whose everywhere-continuous representative is a local martingale for the
  null-augmented filtration, on $[0,T]$ and on all of $\R_{\ge0}$.
\item \emph{It\^o's formula without growth conditions} (\cref{sec:unrestricted}). For every $f$ whose
  six partial derivatives above are continuous, with no bound of any kind,
  $f(t,B_t)-f(0,B_0)-\int_0^t(f_t+\tfrac12f_{xx})(s,B_s)\,\dd s$ is a continuous local martingale, by
  localization at exit times.
\item \emph{It\^o processes} (\cref{sec:itoprocess}). Quadratic variation of an It\^o process with adapted
  volatility, Riemann--Stieltjes sums against an It\^o integral, and It\^o's formula for
  $X=x_0+\int b\,\dd s+\int\sigma\,\dd B$ with bounded adapted continuous coefficients.
\item \emph{The range and the uniqueness of the integral} (\cref{sec:operator}). Martingale
  representation: every square-integrable $\F^B_T$-measurable $F$ is $\E[F]+\ito_T(\varphi)$ for a
  unique $\varphi\in\HT$. Integration against $M=\varphi\bulletB B$ with its bracket. And agreement with
  the \texttt{BrownianMotion} package's axiomatic stochastic integral.
\item \emph{What the calculus carries} (\cref{sec:consumers}). Girsanov's theorem for bounded
  predictable drifts, and existence and uniqueness for Lipschitz stochastic differential equations on
  short horizons.
\end{enumerate}
\Cref{sec:scope} sets the boundary precisely: which parts of the classical theory are present, which are
absent, and why the title says ``a'' calculus rather than ``the'' calculus.

\paragraph{Changes from the second version}
Sections~\ref{sec:unrestricted}--\ref{sec:consumers} are new, as are the growth-condition and process
forms of It\^o's formula in \cref{sec:ito}. The bounded-derivative formula now names its integrand
(the second version displayed an existentially quantified $g_{f'}$ whose prose described more than the
statement said). For the same reason the keystone $\int_0^TB\,\dd B$ is now stated as what the Lean proves,
with the named form derived from the growth-condition formula (\cref{sec:keystone}). The pre-Brownian
interface is \lean{IsPreBrownianReal}, now in Mathlib; the predictable rectangles are consumed from the
upstream package; and the artifact, pins and sizes are those of release \texttt{v1.5.0}. Smaller
descriptions are corrected against the source: the $\sigma$-algebra the rectangles generate is identified
here, not by an upstream lemma; the maximal inequality is Doob's weak-type inequality applied to simple
processes; and the closed form of $\int B\,\dd B$ preceded the general Riemann bridge rather than consuming
it. The scope section is rewritten. Of the limitations it recorded, the restriction of It\^o's formula to
bounded derivatives is lifted; the unrestricted formula, which the second version placed beyond the development in
its $C^2$ form, is now proved for $C^3$ functions in local-martingale form, while the unrestricted $C^2$ formula
for $f(t,B_t)$ remains open.

%======================================================================
\section{Setting and verification discipline}\label{sec:setting}

\subsection{Foundations consumed}
The development is built in Lean~4~\cite{deMouraUllrich2021} over Mathlib~\cite{Mathlib2020} and the
\texttt{BrownianMotion} package~\cite{BrownianMotionRepo,DegenneEtAl2025}, whose construction of Brownian
motion this work consumes rather than reproves; the artifact pins the package at commit
\texttt{0d5b6eb}. The pre-Brownian interface itself now lives in Mathlib. From the package the development
takes the continuous modification of Brownian motion, the simple processes (\lean{SimpleProcess}) and
elementary predictable sets, the predictable rectangles of a filtration and the proof that they form a
semiring of sets, the local-martingale predicate \lean{IsLocalMartingale}, Doob's maximal inequality, and
the axiomatic stochastic integral of \cref{sec:characterisation}. Time is indexed by $\R_{\ge0}$ (Lean's
\lean{NNReal}); Lebesgue measure on time enters through \lean{timeMeasure}, restricted to $(0,T]$.

\subsection{Auditability and reproducibility}
Verification follows the library's standing contract~\cite{CoelhoLibrary}. No \texttt{sorry} or
\texttt{admit} occurs anywhere in the development. Two build-enforced audits run \texttt{\#print axioms}
under \texttt{\#guard\_msgs}: a curated one over headline results, and a generated one that pins every
library constant a benchmark proof cites. Every numbered theorem of this paper, and every result it names
as a headline, is pinned by one of them to \texttt{propext}, \texttt{Classical.choice} and
\texttt{Quot.sound}, and a \texttt{sorry} or extra axiom beneath any of them fails the build. The \texttt{BrownianMotion} package has \texttt{sorry}s of its own at the pin; the
audit is what certifies that none lies beneath a result reported here. A hash ledger records the exact
source inputs and toolchain pins each benchmark entry was last verified under, so a result cannot go
stale silently when a dependency changes. The statements displayed below are transcriptions of the Lean
statements; where a hypothesis package is abbreviated in prose, the Lean declaration is the authority,
and \cref{app:names} maps every displayed result to its declaration, module and benchmark entry.

\subsection{Scale}
Counted as the import closure of its headline modules within the library, less the finance vocabulary
that closure draws in, the It\^o tower is 81 modules and about 27,000 lines of Lean, with some 830
theorems and lemmas; it lies downstream of Mathlib and \texttt{BrownianMotion} and upstream of the
library's pricing layers. The 26 modules of the second version account for about 8,100 of those lines.
\Cref{tab:modules} groups the principal modules by role.

\begin{table}[!t]
\centering\small
\begin{tabularx}{\linewidth}{@{}>{\raggedright\arraybackslash}p{0.21\linewidth}>{\raggedright\arraybackslash}X@{}}
\toprule
Role & Principal modules \\
\midrule
Discrete engine & \file{DiscreteIto}, \file{DiscreteItoPolynomial}, \file{ItoSquaringIdentity} \\
Quadratic variation & \file{QuadraticVariationL2}, \file{WeightedQuadraticVariation},
  \file{BrownianQuadraticVariation}, \file{ItoProcessQV}, \file{AdaptedQuadraticVariation} \\
Integral (isometry) & \file{ItoIntegralL2}, \file{ItoIntegralCLM}, \file{ItoIsometryAdapted},
  \file{ItoIntegralBrownian}, \file{ItoIntegralL2Dense}, \file{ItoIntegralCovariation},
  \file{ItoIntegralLocality} \\
Riemann bridges & \file{ItoIntegralRiemannBridge}, \file{ItoIntegralRiemannBridgeTD},
  \file{ItoIntegralRiemannBridgeAdapted}, \file{AdaptedProcessToLp} \\
Integral as process & \file{ItoIntegralProcess}, \file{ItoIntegralProcessGeneral},
  \file{ItoIntegralProcessMartingale}, \file{ItoIntegralProcessIsometry} \\
It\^o's formula in $L^2$ & \file{ItoFormulaC2}, \file{ItoFormulaCLM}, \file{ItoFormulaRemainder},
  \file{ItoFormulaSquaredL2}, \file{ItoFormulaTD}, \file{ItoFormulaTDRemainder},
  \file{ItoFormulaProcess}, \file{ItoFormulaLocalized}, \file{ItoFormulaGBM} \\
Pathwise regularity & \file{ItoIntegralProcessContinuousModification},
  \file{ItoIntegralProcessLocalMartingale}, \file{ItoIntegralProcessLocalMartingaleGeneral},
  \file{ItoIntegralProcessL2Infinite}, \file{ItoIntegralProcessLocalMartingaleInfinite} \\
Localization & \file{ExitTime}, \file{ItoFormulaUnrestricted}, \file{ItoFormulaUnrestrictedLocMart} \\
It\^o processes & \file{AdaptedStochasticIntegralFreezing}, \file{AdaptedRiemannStieltjes},
  \file{ItoFormulaItoProcess}, \file{ItoFormulaAdapted}, \file{DriftProcessModification} \\
Representation & \file{MartingaleRepresentation}, \file{DoleansStepRepresentation},
  \file{WienerExponentialTotality}, \file{BrownianCylinderGeneration} \\
Integration against $\varphi\bulletB B$ & \file{ItoIntegralAgainstMartingale}, \file{LpMulIsometry},
  \file{PointwiseBracket}, \file{BracketCompensator}, \file{StochasticIntegralCharacterisation} \\
Girsanov & \file{ExpMartingaleQBrownian}, \file{GirsanovConstantTheta}, \file{GirsanovSimpleTheta},
  \file{GirsanovAdaptedTheta}, \file{GirsanovPredictableTheta} \\
SDEs & \file{ItoProcessPredictable}, \file{SDEExistence}, \file{SDEPathwise}, \file{SDEUniqueness} \\
\bottomrule
\end{tabularx}
\caption{The It\^o tower by role: principal modules, all under \texttt{MathFin/Foundations/}. The full
tower, the import closure of these modules less the finance vocabulary, is 81 modules and about 27,000
lines.}
\label{tab:modules}
\end{table}

%======================================================================
\section{The classical theory}\label{sec:classical}

We recall the objects the formalization reproduces, in the form they take in the standard
references~\cite{KaratzasShreve1991,Ito1944}, so that the theorems below can be read against them.

Let $(B_t)_{t\ge0}$ be a Brownian motion on $(\Omega,\F,\mu)$ with natural filtration $(\F_t)$. For a simple
adapted process $\varphi_s=\sum_k\xi_k\ind_{(t_k,t_{k+1}]}(s)$, with $\xi_k$ bounded and
$\F_{t_k}$-measurable, the It\^o integral is the finite sum
$\int_0^T\varphi\,\dd B=\sum_k\xi_k(B_{t_{k+1}}-B_{t_k})$. The defining computation is the \emph{It\^o
isometry}: because each $\xi_k$ is independent of the forward increment and increments over disjoint
intervals are independent,
\[
  \E\Bigl[\Bigl(\sum_k\xi_k(B_{t_{k+1}}-B_{t_k})\Bigr)^2\Bigr]=\sum_k\E[\xi_k^2]\,(t_{k+1}-t_k)
  =\E\int_0^T\varphi_s^2\,\dd s,
\]
the cross terms vanishing by adaptedness and the diagonal contributing $\E[\xi_k^2](t_{k+1}-t_k)$ because
$\E[(B_{t_{k+1}}-B_{t_k})^2\mid\F_{t_k}]=t_{k+1}-t_k$. The map $\varphi\mapsto\int_0^T\varphi\,\dd B$ is
thus an isometry from the simple processes, carrying the $L^2(\lambda\otimes\mu)$ norm on the predictable
side, into $L^2(\mu)$; as the simple processes are dense in the predictable $L^2$ space, it extends to a
linear isometry on all of it (\cref{sec:integral}).

As a function of its upper limit, $t\mapsto\int_0^t\varphi\,\dd B$ is a continuous $L^2$ martingale with
$\E[(\int_0^t\varphi\,\dd B)^2]=\E\int_0^t\varphi_s^2\,\dd s$ (\cref{sec:process,sec:pathwise}). The
quadratic variation of Brownian motion is $[B]_t=t$: along partitions of $[0,t]$ with mesh tending to
zero, $\sum_k(B_{t_{k+1}}-B_{t_k})^2\to t$. It\^o's formula states that for $f\in C^{1,2}$
\[
  f(t,B_t)-f(0,B_0)=\int_0^tf_x(s,B_s)\,\dd B_s+\int_0^t\bigl(f_t+\tfrac12f_{xx}\bigr)(s,B_s)\,\dd s,
\]
the second-order term surviving because $(\dd B)^2=\dd t$ in the mean-square sense made exact by
$[B]_t=t$. When $f_x(\cdot,B)$ is not square-integrable the stochastic integral is a local martingale,
obtained by stopping at the exit times of compact sets (\cref{sec:unrestricted}). Finally, the integral
is onto: every centred square-integrable $\F^B_T$-measurable random variable is
$\int_0^T\varphi\,\dd B$ for a unique predictable $\varphi$, the martingale representation theorem
(\cref{sec:operator}).

%======================================================================
\section{Design of the formalization}\label{sec:design}

Four representation choices shape the development and are worth stating before the theorems that rest on
them.

\paragraph{The integrand, after the upstream package}
The simple integrands are not re-encoded here. \texttt{BrownianMotion} provides \lean{SimpleProcess}, a
bounded value at time zero together with finitely many bounded values on intervals, each measurable at the
interval's left endpoint, and the predictable rectangles of a filtration together with the proof that they
form a semiring of sets. This work consumes both: the elementary It\^o integral is defined on
\lean{SimpleProcess}, the rectangle family is the upstream one, and the $\pi$-system property needed for
Dynkin's theorem is projected from the upstream semiring. What is built here is the layer above: the
identification of the $\sigma$-algebra the rectangles generate, the orthogonal-complement density and the
isometric extension.

\paragraph{Time as \texttt{NNReal}}
The time axis is $\R_{\ge0}$, not $\R$ with a positivity side condition. Increments $B_t-B_s$ then carry
$s\le t$ in their type-level data rather than as a hypothesis to discharge, and the bounded horizon
enters through a single finite measure (\lean{timeMeasure_T}, of total mass $T$) rather than through
interval hypotheses threaded across lemmas.

\paragraph{A trimmed measure}
The integrand space is the $L^2$ space of a measure that has already been restricted to the predictable
$\sigma$-algebra, rather than an ambient $L^2$ space with a sub-$\sigma$-algebra measurability side
condition. Predictability is then a property of membership in $\HT$, not a hypothesis threaded through
every limit: a Cauchy sequence in $\HT$ has its limit in $\HT$, so when It\^o's formula realizes
$s\mapsto f'(B_s)$ as an $L^2$ limit of step processes the limit's predictability is free. The cost is a
one-time investment in the trimmed-measure $L^2$ API; the saving recurs at every density and convergence
argument downstream.

\paragraph{One construction pattern, used throughout}
The integral, the integral as a process and the integral against an It\^o integral are built by the same
idiom: a norm-preserving linear map on a dense subspace, then its unique continuous linear extension
(Mathlib's \lean{LinearMap.extendOfNorm}). Every identity proved on simple processes then transfers to
all of $\HT$ by continuity and density, through the equalizer property of dense ranges
(\lean{DenseRange.equalizer}). That mechanism is what turns \cref{sec:process} into a sequence of short
arguments.

%======================================================================
\section{The \texorpdfstring{$L^2$}{L2} It\^o integral}\label{sec:integral}

\subsection{The integrand space}
Fix $T\in\R_{\ge0}$ and write $\lambda_T$ for Lebesgue measure on time restricted to $(0,T]$. The
predictable rectangles
\[
  \mathcal R=\bigl\{(s,t]\times A:\ s<t,\ A\in\F_s\bigr\}\cup\bigl\{\{0\}\times A_0:\ A_0\in\F_0\bigr\}
\]
of the natural filtration form a $\pi$-system (\lean{isPiSystem_predictableRect}) that generates the
predictable $\sigma$-algebra $\mathcal P$ (\lean{generateFrom_predictableRect}, proved here directly);
\cref{lst:rect} shows the definition, which is now the upstream one. The integrand space is
\[
  \HT=L^2\bigl((0,T]\times\Omega,\ \mathcal P,\ (\lambda_T\otimes\mu)\restriction_{\mathcal P}\bigr),
\]
realized as the $L^2$ space of the product measure trimmed to $\mathcal P$ (\lean{trimMeasure_T}), and
the It\^o integral is a continuous linear map out of it (\cref{lst:clm}).

\begin{leancode}[caption={The predictable rectangles, consumed from \texttt{BrownianMotion} (\texttt{Foundations/ItoIntegralCLM}).},label=lst:rect]
def predictableRect (hBmeas : ∀ t, Measurable (B t)) :
    Set (Set (ℝ≥0 × Ω)) :=
  (ItoIntegralL2.natFiltration (mΩ := mΩ) hBmeas).predictableRectangles

lemma isPiSystem_predictableRect (hBmeas : ∀ t, Measurable (B t)) :
    IsPiSystem (predictableRect (mΩ := mΩ) hBmeas) :=
  (MeasureTheory.Filtration.isSetSemiring_predictableRectangles _).isPiSystem
\end{leancode}

\begin{leancode}[caption={The It\^o integral and its isometry (\texttt{Foundations/ItoIntegralCLM}).},label=lst:clm]
noncomputable def itoIntegralCLM_T (hB : IsPreBrownianReal B μ) (T : ℝ≥0)
    (hBmeas : ∀ t, Measurable (B t)) :
    Lp ℝ 2 (trimMeasure_T (μ := μ) T hBmeas) →L[ℝ] Lp ℝ 2 μ :=
  (itoAssembly_T hB T hBmeas).extendOfNorm (simpleAssembly_T (μ := μ) T hBmeas)

theorem itoIntegralCLM_T_norm (hB : IsPreBrownianReal B μ) (T : ℝ≥0)
    (hBmeas : ∀ t, Measurable (B t))
    (f : Lp ℝ 2 (trimMeasure_T (μ := μ) T hBmeas)) :
    ‖itoIntegralCLM_T hB T hBmeas f‖ = ‖f‖
\end{leancode}

\subsection{Simple processes, density, isometry}
A simple adapted process is a finite combination of indicator rectangles $\xi_s\cdot\ind_{(s,t]}$ with
$\xi_s$ bounded and $\F_s$-measurable; \lean{TBoundedSP} is the submodule of those supported in $[0,T]$ and
\lean{simpleProcessL2_T} sends it into $\HT$. The elementary integral maps such a process to
$\sum\xi_s(B_t-B_s)$ (\lean{itoAssembly_T}); independence and the Gaussian increment moments give the
elementary It\^o isometry (\lean{assembly_isometry_T}).

Density of the simple processes in $\HT$ (\lean{simpleAssembly_T_denseRange}) is the technical core of
the construction and the part of it not inherited from the upstream package. It is proved by the
orthogonal-complement method. Take $\psi\in\HT$ orthogonal to every elementary integrand. Testing against
the indicator of a rectangle $(a,b]\times F$ shows that the integral of $\psi$ over that rectangle
vanishes; the measurable sets on which $\int\psi=0$ form a $\lambda$-system containing the rectangle
$\pi$-system, so Dynkin's theorem extends the vanishing to all of $\mathcal P$
(\lean{setIntegral_eq_zero_of_orthogonal_pred}), forcing $\psi=0$ almost everywhere. The It\^o integral
$\ito_T$ (\lean{itoIntegralCLM_T}) is then the unique isometric continuous linear extension of the
elementary integral. The same density, run over the $\sigma$-finite time axis by exhaustion, gives the
unbounded-horizon isometry on $L^2(\R_{\ge0}\times\Omega)$ (\lean{itoIntegralL2_norm}).

\subsection{The keystone example}\label{sec:keystone}
The integral's first consumer was the closed form of $\int_0^TB\,\dd B$: there is an integrand $g_B\in\HT$,
built as the $\HT$-limit of truncated left-endpoint step processes of $s\mapsto B_s$, with
\[
  \ito_T(g_B)=\tfrac12\bigl(B_T^2-B_0^2-T\bigr)\qquad(\text{\lean{itoIntegralCLM_T_brownian}}).
\]
The It\^o correction $-T/2$ emerges as follows: the summation-by-parts identity
\lean{discrete_squaring_identity} writes
\[
  \sum_kB_{s_k}(B_{s_{k+1}}-B_{s_k})=\tfrac12\Bigl(B_T^2-B_0^2-\sum_k(B_{s_{k+1}}-B_{s_k})^2\Bigr),
\]
whose left side is the discretized $\int_0^TB\,\dd B$ and whose squared-increment sum converges in $L^2$ to
$T$ (\cref{sec:qv}); what is left is $\frac12(B_T^2-B_0^2-T)$. Because the left-endpoint coefficients
$B_{s_k}$ are unbounded Gaussians, the sums are not elementary integrals of simple processes; the proof
truncates them to $[-M,M]$ and removes the truncation in the $L^2$ limit, which is how the It\^o integral is
defined for unbounded integrands in any case. The general Riemann bridge, which identifies $\ito_T$ of a
step-process realization of a bounded continuous function of $B$ with the classical partial sums and feeds
It\^o's formula, came afterwards and builds on this file.

The statement records the integral, not the integrand. The construction makes $g_B$ the class of $B$, but
the theorem does not say so, and by martingale representation (\cref{thm:mrt}) the bare existence of
\emph{some} integrand would hold for any centred square-integrable $\F^B_T$-measurable target. The named
form, $\ito_T(B)=\frac12(B_T^2-B_0^2-T)$ with $B$ the class of $(s,\omega)\mapsto B_s(\omega)$, is the instance
$f(t,x)=x^2$ of the growth-condition formula (\cref{thm:ito-growth}), whose statement names its integrand:
it gives $g$ equal to $2B$ almost everywhere with $B_T^2-B_0^2=\ito_T(g)+T$. That instance is a few lines
from \cref{thm:ito-growth} but is not stated separately in the library. The identity remains the smallest
statement in which the calculus is visibly non-classical, and it anchors the remainder analysis of
\cref{sec:ito}.

\subsection{Design note: two isometries}
The library separately maintains a Wiener-integral isometry for deterministic integrands. The two are
deliberately distinct: the Wiener integral needs no filtration and serves the Gaussian-process and
distributional layers, while $\ito_T$ is the adapted object the martingale theory of \cref{sec:process}
requires. Collapsing them would shorten the code and damage the mathematics; the duplication is the
design, and it is recorded as such in the library's dependency blueprint.

%======================================================================
\section{The integral as a process}\label{sec:process}

\subsection{The simple layer}
For a simple process $V$ and $t\le T$, the partially summed elementary integral $(V\bulletB B)_t$ is
adapted (\lean{itoSimpleProcess_adaptedAt}) and a martingale (\lean{itoSimpleProcess_isMartingale}). The
martingale step is the conditional vanishing of adapted-weighted future increments: by the characterizing
property of conditional expectation it suffices that $\int_A\xi(B_{t_1}-B_{t_0})\,\dd\mu=0$ for every
$A\in\F_{t_0}$, and $\xi\ind_A$ is again bounded and $\F_{t_0}$-measurable, so the increment, independent of
$\F_{t_0}$ with mean zero, factors the integral as $\E[\xi\ind_A]\,\E[B_{t_1}-B_{t_0}]=0$. The martingale proof
runs this argument set by set; the conditional statement it amounts to is recorded on its own in
\cref{lst:condexp}. The simple layer also carries the per-time isometry
$\E[(V\bulletB B)_t^2]=\int_0^t\E[V_s^2]\,\dd s$ (\lean{itoSimpleProcess_isometry_time}) and
$L^2$-continuity in $t$ (\lean{itoSimpleProcessLp_l2_continuous}); the next subsection lifts all three to
general integrands at once.

\begin{leancode}[caption={Adapted weights kill future increments (\texttt{Foundations/ItoIntegralProcessMartingale}).},label=lst:condexp]
theorem condExp_adapted_mul_increment (hBmeas : ∀ t, Measurable (B t))
    {t₀ t₁ : ℝ≥0} (ht : t₀ ≤ t₁) {φ : Ω → ℝ}
    (hφ : ItoIsometryAdapted.AdaptedAt B t₀ φ) {C : ℝ} (hC : ∀ ω, |φ ω| ≤ C) :
    μ[fun ω ↦ φ ω * (B t₁ ω - B t₀ ω) | ItoIntegralL2.natFiltration hBmeas t₀]
      =ᵐ[μ] 0
\end{leancode}

\subsection{The key identity}
For general $\varphi\in\HT$ the process is a continuous linear map (\lean{itoProcessCLM}, the extension of
the simple layer), and the development's central structural fact is the following.

\begin{theorem}[\lean{itoProcessCLM_eq_condExpL2}]\label{thm:key}
For every $t\ge0$ and $\varphi\in\HT$, as elements of $L^2(\Omega,\mu)$,
\[
  (\varphi\bulletB B)_t=\E\bigl[\ito_T(\varphi)\bigm|\F_t\bigr],
\]
the right side realized by Mathlib's $L^2$ conditional-expectation projection \lean{condExpL2}.
\end{theorem}

\begin{leancode}[caption={The integral process is the conditional projection of its terminal value (\texttt{Foundations/ItoIntegralProcessGeneral}).},label=lst:key]
theorem itoProcessCLM_eq_condExpL2 (T t : ℝ≥0) (hBmeas : ∀ u, Measurable (B u))
    (φ : Lp ℝ 2 (trimMeasure_T (μ := μ) T hBmeas)) :
    itoProcessCLM hB T t hBmeas φ
      = (condExpL2 ℝ ℝ ((natFiltration hBmeas).le t)
          (itoIntegralCLM_T hB T hBmeas φ) : Lp ℝ 2 μ)
\end{leancode}

The proof is two lines of mathematics: both sides are continuous linear in $\varphi$; on the dense simple
subspace they agree, because there the identity is the simple layer's martingale property; conclude by the
equalizer property of dense ranges. Everything else about the process is a corollary of
\cref{thm:key} and standard properties of conditional expectation:
\begin{itemize}
\item \emph{adaptedness}, since a \lean{condExpL2} projection lands in the $\F_t$-measurable subspace;
\item the \emph{martingale property} (\lean{itoIntegralProcessGen_isMartingale}), by the tower property;
\item the \emph{contraction bound} $\norm{(\varphi\bulletB B)_t}_{L^2}\le\norm{\varphi}_{\HT}$
  (\lean{itoProcessCLM_norm_le}), since conditional expectation is a contraction;
\item the \emph{terminal isometry} at $t=T$ (\lean{itoProcessCLM_norm_terminal});
\item \emph{$L^2$-continuity} of $t\mapsto(\varphi\bulletB B)_t$ (\lean{itoIntegralProcessGen_l2_continuous}),
  by a $t$-uniform contraction estimate against an approximating simple sequence.
\end{itemize}
Defining the process as a projection of its terminal value is what turns the martingale theory of the
It\^o integral into a sequence of corollaries rather than a sequence of constructions. Two further
structural facts follow the same pattern. The isometry polarizes:
$\E[\ito_T(\varphi)\,\ito_T(\psi)]=\int\varphi\psi\,\dd(\lambda_T\otimes\mu)$, the covariation of two It\^o
integrals (\file{ItoIntegralCovariation}). And the integral is local in time: if $\varphi$ vanishes on
$[0,a]$ and $Z$ is $\F_a$-measurable, then $\ito_T(Z\varphi)=Z\,\ito_T(\varphi)$ (\file{ItoIntegralLocality}),
the ``freeze the past, integrate the future'' step behind the Dol\'eans exponentials of
\cref{sec:representation}.

\subsection{The time-indexed It\^o isometry}
The terminal isometry is the classical $\norm{\ito_T(\varphi)}^2_{L^2}=\norm{\varphi}^2_{\HT}$. Its
refinement to intermediate times, the energy law of the integral as a process, is the second fact every
textbook calls ``the It\^o isometry''.

\begin{theorem}[\lean{itoProcessCLM_norm_sq}]
For every $t\le T$ and $\varphi\in\HT$,
\[
  \E\bigl[(\varphi\bulletB B)_t^2\bigr]=\int_{(0,t]\times\Omega}\varphi^2\,\dd(\lambda_T\otimes\mu)\restriction_{\mathcal P}
  =\int_0^t\E[\varphi_s^2]\,\dd s .
\]
\end{theorem}

\begin{leancode}[caption={The time-indexed It\^o isometry (\texttt{Foundations/ItoIntegralProcessIsometry}).},label=lst:timeiso]
theorem itoProcessCLM_norm_sq {t T : ℝ≥0} (htT : t ≤ T) (hBmeas : ∀ u, Measurable (B u))
    (φ : Lp ℝ 2 (trimMeasure_T (μ := μ) T hBmeas)) :
    ‖itoProcessCLM hB T t hBmeas φ‖ ^ 2
      = ∫ z in (Set.Ioc 0 t ×ˢ (Set.univ : Set Ω)), (φ z) ^ 2
          ∂(trimMeasure_T (μ := μ) T hBmeas)
\end{leancode}

The proof is again the equalizer pattern, against a different pair of continuous functionals:
$\varphi\mapsto\norm{(\varphi\bulletB B)_t}^2$ and $\varphi\mapsto\int_{(0,t]}\varphi^2$ agree on simple
processes by a band reconciliation that splits the trimmed measure across $(0,t]\cup(t,T]$
(\lean{itoSimpleProcessLp_band_isometry}, \lean{integral_rectTerm_mul_band}). Stating it for general
integrands at every $t$, not only at $t=T$ and not only for simple processes, is what makes the integral
as a process a complete object rather than a martingale with a missing variance law.

%======================================================================
\section{It\^o's formula in \texorpdfstring{$L^2$}{L2}}\label{sec:ito}

\subsection{Architecture}
The route to It\^o's formula is discrete-to-continuous (\cref{fig:arch}).
\begin{enumerate}
\item \emph{Discrete identity.} For any partition, second-order Taylor expansion of $f$ along the
  increments gives an exact discrete It\^o identity (\file{DiscreteIto}) with an explicit Lagrange
  remainder.
\item \emph{Weighted quadratic variation.} The term $\sum_kf''(B_{t_k})(\Delta_kB)^2$ converges in $L^2$
  to $\int_0^Tf''(B_s)\,\dd s$; this is proved for adapted weight processes
  (\file{WeightedQuadraticVariation}), reusing the Gaussian moment estimates of the unweighted limit.
\item \emph{Remainder.} The third-order remainder is $O(1/n)$ in $L^2$ under a bound on $f'''$
  (\file{ItoFormulaRemainder}); the time-dependent variant needs a two-dimensional remainder with mixed
  partials, also $O(1/n)$ (\file{ItoFormulaTDRemainder}).
\item \emph{Riemann bridge.} The discretized stochastic term is $\ito_T$ applied to a step-process
  realization of $f'(B)$.
\item \emph{Limit.} The step integrands converge in $\HT$ to the class of $s\mapsto f'(B_s)$. Predictability
  of the limit is free (the limit is taken inside $\HT$); its existence for bounded continuous data is
  \lean{itoIntegralCLM_T_of_bdd_cont}, a dominated-convergence argument on the trimmed measure.
\end{enumerate}

\begin{figure}[tbp]
\centering
\begin{tikzpicture}[
  box/.style={draw=mfaccent!70, fill=mfbg, rounded corners=2pt, align=center,
    inner sep=3pt, font=\footnotesize, minimum height=9mm, text width=26mm,
    execute at begin node={\hyphenpenalty=10000\exhyphenpenalty=10000}},
  wide/.style={box, text width=29mm},
  arr/.style={-{Stealth[length=2mm]}, draw=mfaccent}]
\node[box]  (pi)   at (0,0)      {predictable $\pi$-system and density};
\node[box]  (int)  at (3.75,0)    {$L^2$ It\^o integral $\ito_T$};
\node[box]  (proc) at (7.5,0)    {integral as a process};
\node[wide] (path) at (11.4,0)   {continuous local-martingale modification};
\node[box]  (ito)  at (3.75,-2.1) {It\^o's formula in $L^2$};
\node[wide] (loc)  at (11.4,-2.1) {exit-time localization: no growth condition};
\node[box]  (disc) at (0,-4.1)   {discrete It\^o identity};
\node[box]  (qv)   at (2.95,-4.1) {QV in $L^2$ ($2t^2/n$)};
\node[box]  (wqv)  at (5.9,-4.1) {weighted QV in $L^2$};
\node[box]  (rem)  at (8.85,-4.1) {remainder $O(1/n)$};
\draw[arr] (pi) -- (int);
\draw[arr] (int) -- (proc);
\draw[arr] (proc) -- (path);
\draw[arr] (int) -- (ito);
\draw[arr] (qv) -- (wqv);
\draw[arr] (disc.north) -- (ito.south west);
\draw[arr] (wqv.north) -- (ito.south);
\draw[arr] (rem.north) -- (ito.south east);
\draw[arr] (ito) -- (loc);
\draw[arr] (path) -- (loc);
\end{tikzpicture}
\caption{Proof architecture. The isometry layer (top) builds the integral, the process and its continuous
modification; the discrete-to-continuous layer (bottom) supplies the vanishing terms, which meet at It\^o's
formula; localization combines the formula with the pathwise process to remove every growth condition.}
\label{fig:arch}
\end{figure}

\subsection{The fluctuation engine: quadratic variation in \texorpdfstring{$L^2$}{L2}}\label{sec:qv}
The second-order term of It\^o's formula is deterministic in the limit for one reason: the squared
increments of Brownian motion have vanishing fluctuation in $L^2$. Along the uniform partition of $[0,t]$
into $n$ pieces of width $\Delta=t/n$,
\[
  \Bigl\lVert\sum_k(B_{s_{k+1}}-B_{s_k})^2-t\Bigr\rVert_{L^2(\mu)}^2=\frac{2t^2}{n}\longrightarrow0
  \qquad(\text{\lean{tendsto_qv}}).
\]
The value is exact, not asymptotic: for any partition $0=s_0\le\dots\le s_n$ the mean-square error is
$\sum_k2(s_{k+1}-s_k)^2$ (\lean{sum_increment_sq_sub_sq_integral}), which is $2t^2/n$ for the uniform one. Writing $Y_k=(B_{s_{k+1}}-B_{s_k})^2-\Delta$ for the centred squared
increment, the $Y_k$ over disjoint intervals are independent, so the cross terms vanish and
$\E[(\sum_kY_k)^2]=\sum_k\E[Y_k^2]$. Each term is the Gaussian kurtosis,
$\E[Y_k^2]=\E[(B_{s_{k+1}}-B_{s_k})^4]-\Delta^2=3\Delta^2-\Delta^2=2\Delta^2$, using $\E[X^4]=3(\operatorname{Var}X)^2$ for
a centred Gaussian (\lean{integral_pow4_gaussianReal}, transported to the increment law by
\lean{integral_increment_pow4}). Summing, $n\cdot2\Delta^2=2t^2/n$. The fourth-moment constant $3$, and nothing
else, is the reason the quadratic variation is $t$ rather than $0$; this is the formal content of
$(\dd B)^2=\dd t$.

\subsection{The weighted second-order term}
It\^o's formula needs this limit carrying a weight: for a bounded adapted process $w$ with continuous
paths,
\[
  \sum_kw_{t_k}(B_{t_{k+1}}-B_{t_k})^2\longrightarrow\int_0^Tw_s\,\dd s\quad\text{in }L^2
  \qquad(\text{\lean{tendsto_weighted_qv_process}}),
\]
instantiated at $w_s=f''(B_s)$ for the autonomous formula and $w_s=f_{xx}(s,B_s)$ for the time-dependent
one. The proof splits the difference into a fluctuation term and a Riemann term. The fluctuation term
$\sum_kw_{t_k}((B_{t_{k+1}}-B_{t_k})^2-\Delta_k)$ vanishes in $L^2$ by the orthogonality-plus-kurtosis estimate
above, now carrying the $\F_{t_k}$-measurable bounded weight: adaptedness of $w_{t_k}$ is exactly what keeps the
cross terms zero (\lean{integral_adapted_mul_centered_sq}). The Riemann term vanishes pathwise by continuity
of $s\mapsto w_s(\omega)$, then in $L^2$ by dominated convergence against the bound on $w$. Neither half uses
where the weight came from, only that it is adapted, bounded and path-continuous, which is why one lemma
serves the autonomous, the time-dependent, and later the It\^o-process formula.

\subsection{The remainder}
What survives after the first- and second-order terms is a sum of cubic Taylor remainders, and it must vanish
in $L^2$. The per-step remainder $R_k$ obeys $\abs{R_k}\le C_3\abs{B_{s_{k+1}}-B_{s_k}}^3$ under
$\abs{f'''}\le C_3$ (\lean{abs_discreteTaylorRemainder_le}). Hence $\E[R_k^2]=O(\E[(\Delta_kB)^6])=O(\Delta^3)$ by the
Gaussian sixth moment (\lean{integral_pow6_gaussianReal}), and Cauchy--Schwarz across the $n$ terms gives
\[
  \E\Bigl[\Bigl(\sum_kR_k\Bigr)^2\Bigr]\le n\sum_k\E[R_k^2]=O\bigl(n\cdot n\cdot(T/n)^3\bigr)=O(1/n)\longrightarrow0 .
\]
The third order is the only place a third derivative is used. The time-dependent variant runs the same
estimate on a two-dimensional Taylor remainder with mixed partials.

\subsection{The bounded-derivative formula}

\begin{theorem}[\lean{ito_formula_L2_bddDeriv}]\label{thm:ito-bdd}
Let $B$ have measurable coordinates and continuous paths, and let $f\colon\R\to\R$ be three times
differentiable with $\abs{f'}\le C_1$, $\abs{f''}\le C_2$, $\abs{f'''}\le C_3$. Then there is
$g\in\HT$ with $g=f'(B_s(\omega))$ for $(\lambda_T\otimes\mu)\restriction_{\mathcal P}$-almost every $(s,\omega)$,
such that $\mu$-almost surely
\[
  f(B_T)-f(B_0)=\ito_T(g)+\tfrac12\int_{(0,T]}f''(B_s)\,\dd s .
\]
\end{theorem}

\begin{leancode}[caption={It\^o's formula, bounded-derivative class (\texttt{Foundations/ItoFormulaCLM}).},label=lst:itobdd]
theorem ito_formula_L2_bddDeriv
    (hBmeas : ∀ t, Measurable (B t)) (hBcont : ∀ ω, Continuous (fun s : ℝ≥0 ↦ B s ω))
    (T : ℝ≥0) {f f' f'' f''' : ℝ → ℝ}
    (hf : ∀ x, HasDerivAt f (f' x) x) (hf' : ∀ x, HasDerivAt f' (f'' x) x)
    (hf'' : ∀ x, HasDerivAt f'' (f''' x) x)
    {C1 : ℝ} (hf1 : ∀ x, |f' x| ≤ C1) {C2 : ℝ} (hf2 : ∀ x, |f'' x| ≤ C2)
    {C3 : ℝ} (hf3 : ∀ x, |f''' x| ≤ C3) :
    ∃ gf' : Lp ℝ 2 (trimMeasure_T (μ := μ) T hBmeas),
      (⇑gf' =ᵐ[trimMeasure_T (μ := μ) T hBmeas] fun z ↦ f' (B z.1 z.2)) ∧
      (fun ω ↦ f (B T ω) - f (B 0 ω)) =ᵐ[μ]
        (fun ω ↦ (itoIntegralCLM_T hB T hBmeas gf') ω
          + (1 / 2) * ∫ s in Set.Ioc 0 T, f'' (B s ω) ∂ItoIntegralL2.timeMeasure)
\end{leancode}

The first conjunct is what makes the second It\^o's \emph{formula} rather than a decomposition into some
It\^o integral: the stochastic term is the integral of the class of $f'(B)$, not of an unnamed element of
$\HT$.

\begin{proof}[Proof sketch]
Fix the uniform partition of $[0,T]$ into $n$ pieces and expand $f$ to second order along each increment.
The telescoping sum gives the exact discrete identity $f(B_T)-f(B_0)=S_n+\frac12Q_n+R_n$, with
$S_n=\sum_kf'(B_{s_k})(B_{s_{k+1}}-B_{s_k})$, $Q_n=\sum_kf''(B_{s_k})(B_{s_{k+1}}-B_{s_k})^2$ and $R_n$ the cubic
remainder. Each term has a limit: $S_n\to\ito_T(g)$ in $L^2$ because the step integrands converge to $g$ in
$\HT$ and $\ito_T$ is isometric; $Q_n\to\int_0^Tf''(B_s)\,\dd s$ by the weighted quadratic variation; and
$R_n\to0$ by the remainder estimate. The elementary inequality $(a+b+c)^2\le3a^2+3b^2+3c^2$ sends the total
error to zero in $L^2$, so the two sides agree in $L^2$, hence almost surely.
\end{proof}

Path continuity is supplied by the package's continuous modification and is genuinely used: the Riemann
limits run along refining partitions evaluated on paths.

\subsection{The time-dependent formula, and the formula as a process}

\begin{theorem}[\lean{ito_formula_td_L2_bddDeriv}]\label{thm:ito-td}
Let $f\colon\R\times\R\to\R$ have everywhere partial derivatives $f_t,f_x,f_{xx},f_{tt},f_{tx},f_{xxx}$, all
bounded, with $f_x$ and $f_{xx}$ jointly continuous. Then there is $g\in\HT$ equal to
$f_x(s,B_s(\omega))$ almost everywhere such that, $\mu$-almost surely,
\[
  f(T,B_T)-f(0,B_0)=\ito_T(g)+\int_{(0,T]}\bigl(f_t+\tfrac12f_{xx}\bigr)(s,B_s)\,\dd s .
\]
\end{theorem}

Two formalization points. First, joint continuity of $f_t$ is not assumed: it is \emph{derived} from the
bounded mixed partials (\lean{continuous_uncurry_of_bdd_partials}), keeping the hypothesis package at what a
user can check. Second, the time-direction Taylor terms are handled by a boundary-cancellation decomposition
in which three telescoping error sums vanish separately, glued by the three-term Cauchy--Schwarz estimate.

The formula also holds as a process. Under the same hypotheses there is an integrand $F$ on the
half-line such that, for every $t\le T$, $f(t,B_t)-f(0,B_0)=(F\bulletB B)_t+\int_0^t(f_t+\frac12f_{xx})(s,B_s)\,\dd s$
almost surely, the integral process being a martingale with a continuous local-martingale modification
(\lean{ito_formula_td_process}). This statement, unlike \cref{thm:ito-td}, does not name $F$.

\subsection{Exponential growth}\label{sec:growth}
Bounded derivatives exclude the textbook's first examples: $x^2$ has an unbounded first derivative, and the
geometric Brownian motion value function $S_0\exp((r-\sigma^2/2)t+\sigma x)$ has all its derivatives
proportional to $e^{\sigma x}$. The formula is lifted to them by a cutoff localization that reuses the
bounded engine.

\begin{theorem}[\lean{ito_formula_td_localized}]\label{thm:ito-growth}
The conclusion of \cref{thm:ito-td} holds when the six partial derivatives, instead of being bounded,
satisfy $\abs{f_\bullet(t,x)}\le Ce^{\lambda\abs{x}}$ for some $C$ and $\lambda\ge0$, and $f_t$, $f_x$, $f_{xx}$
are jointly continuous.
\end{theorem}

A smooth truncation $\phi_n$, equal to the identity on $[-n,n]$ with bounded derivatives, makes
$f_n(t,x)=f(t,\phi_n(x))$ satisfy the bounded hypotheses, so \cref{thm:ito-td} applies to each $f_n$, with
integrands $g_n$. The boundary and drift terms converge in $L^2(\mu)$ by dominated convergence, the
exponential dominators being integrable because Brownian marginals have all exponential moments. Hence the
integrals $\ito_T(g_n)$ converge and are Cauchy; the isometry transfers Cauchy-ness back to the $g_n$, which
converge in the complete space $\HT$ to some $g$; continuity of $\ito_T$ identifies the limit, and the
almost-everywhere identification of $g$ with $f_x(s,B_s)$ passes to the limit along a subsequence. The
truncation is the antiderivative of a smooth bump, so every derivative bound comes from Mathlib rather than
from explicit calculus. Applied to the exponential, the theorem gives the It\^o decomposition of geometric
Brownian motion (\lean{ito_formula_gbm}) and, for the discounted price, the representation
$e^{-rt}S_t=S_0+\ito_t(\sigma e^{-r\cdot}S)$ with its integrand named (\lean{discountedGBM_eq_itoIntegral}),
which is where the library's Black--Scholes layer meets the It\^o integral.

\subsection{Supporting results}
\paragraph{Quadratic variation}
For a process with measurable coordinates and centred Gaussian increments of variance $t-s$ (neither
independence of increments nor path continuity is assumed), the expected squared-increment sums along
equipartitions of $[0,t]$ converge to $t$ (\lean{qv_equals_t}); the $L^2$ form along uniform partitions is
\lean{tendsto_qv}, and convergence in probability is \lean{tendstoInMeasure_qv}. The $L^1$ statement needs only the second moment of the increments, and so strictly
strengthens the classical claim.

\paragraph{Expectation-level It\^o}
For $f$ with $f$, $f'$ and $f''$ bounded and $f''$ continuous,
$\E[f(B_t)]=f(0)+\frac12\int_0^t\E[f''(B_s)]\,\dd s$ for $t>0$ (\lean{expectation_ito}). It is the bridge between the stochastic tower of this paper and the heat-kernel
tower of the parent library, where the Black--Scholes equation is derived through Feynman--Kac.

\subsection{Notes on the proof engineering}
The density argument is the only place a monotone-class induction was unavoidable in the construction of the
integral; expressing it against a trimmed measure, so that the inductively built sets live in the right
$\sigma$-algebra by construction, is what kept it tractable. The dominated-convergence step that produces the
integrand of It\^o's formula is delicate because the dominating function must itself be predictable; the
trimmed-measure design is what makes the domination land inside $\HT$ without a separate measurability
obligation.

%======================================================================
\section{From \texorpdfstring{$L^2$}{L2} to pathwise: the continuous local martingale}\label{sec:pathwise}

\Cref{sec:process,sec:ito} are an $L^2$ theory: the process $t\mapsto(\varphi\bulletB B)_t$ is continuous
\emph{as a curve in $L^2(\mu)$}, and its laws are identities of $L^2$ elements. The classical object is
sharper: for almost every $\omega$ it has a continuous path, and it is a continuous local martingale in the
pathwise sense. This section builds that object, first on $[0,T]$ and then on the half-line.

\subsection{The pathwise gap and the maximal inequality}
$L^2$-continuity in $t$ does not by itself give a continuous path: it controls each time separately, not the
trajectory. The bridge is a maximal inequality. \texttt{BrownianMotion} supplies Doob's weak-type maximal
inequality for right-continuous martingales; applied to the simple integral processes, with the $L^1$ norm
bounded by the $L^2$ norm and the isometry, it gives
$\mu\{\sup_{t\le T}\abs{(V\bulletB B)_t}\ge\varepsilon\}\le\varepsilon^{-1}\norm{V}_{\HT}$
(\lean{itoSimpleProcess_maximal_prob}). A geometric subsequence makes these tails summable, and Borel--Cantelli
makes the approximating step-process integrals converge uniformly on $[0,T]$ for almost every $\omega$. A uniform limit of continuous functions is continuous, so the limit path is
continuous off a null set.

\subsection{The continuous modification}

\begin{theorem}[\lean{exists_continuous_modification_itoProcess}]
Let $B$ have every path continuous. For every $\varphi\in\HT$ there is a process $\widetilde X\colon\R_{\ge0}\times\Omega\to\R$ with $\mu$-almost
surely continuous paths such that $\widetilde X_t=(\varphi\bulletB B)_t$ in $L^2(\mu)$ for every $t\le T$.
\end{theorem}

The representative \lean{itoContinuousMod} is the pointwise limit of the approximating step integrals on the
good set where they converge uniformly. It equals $(\varphi\bulletB B)_t$ almost surely at each $t$ by uniqueness of
limits in measure, and its paths are continuous by the uniform-convergence theorem; the supporting bound on the
running supremum over $[0,t]$ is \lean{norm_le_iSup_Iic}.

\subsection{The local martingale on the null-augmented filtration}
A modification fixes the paths but can break adaptedness, since the continuous representative is defined
through a limit taken off a null set. The remedy is to complete the filtration: adjoining the $\mu$-null sets
to each $\F_t$ gives $\F^\mu_t=\F_t\vee\mathcal N$ (\lean{augFiltration}). Completeness is all that is added;
right-continuity, the other half of the usual conditions, is neither added nor needed.

\begin{theorem}[\lean{exists_continuous_localMartingale_modification}]
Let $B$ have every path continuous. For every $\varphi\in\HT$ there is a process $X$, adapted to
$(\F^\mu_t)$ and with everywhere-continuous paths, such that $X_t=(\varphi\bulletB B)_t$ in $L^2(\mu)$ for each
$t\le T$, and $X$ is a local martingale for $(\F^\mu_t)$.
\end{theorem}

Two points fix the content. The underlying $L^2$ process is a \emph{true} martingale (\cref{sec:process}); what
the theorem adds is an everywhere-continuous representative, in the form the local-martingale interface of
\texttt{BrownianMotion} takes, whose requirement of paths continuous for \emph{every} $\omega$ is what forces an
everywhere-defined representative. And the martingale property survives the completion, because a conditional
expectation is unchanged when the filtration is enlarged by null sets (\lean{condExp_sup_nulls}).

\subsection{The unbounded horizon}
The bounded-horizon representatives are consistent: for $m\le n$ the horizon-$m$ and horizon-$n$ continuous
processes agree, almost surely and pathwise, on $[0,m)$, because both are continuous and both modify the same
$L^2$ process there. Gluing them gives a single process on the half-line, reading at time $t$ the horizon
$\lceil t\rceil+1$.

\begin{theorem}[\lean{exists_continuous_localMartingale_modification_infinite}]
For every $\varphi$ in the unbounded-horizon integrand space there is a process $X$ with everywhere-continuous
paths such that $X_t=(\varphi\bulletB B)_t$ in $L^2(\mu)$ for every $t$, and $X$ is a local martingale for
$(\F^\mu_t)$ on $\R_{\ge0}$.
\end{theorem}

This is the It\^o integral in the form the classical theory gives it: a continuous local martingale on the
half-line. The glued process is in fact a true martingale with continuous paths
(\lean{exists_continuous_martingale_modification_infinite}), which is what the localization of
\cref{sec:unrestricted} uses.

%======================================================================
\section{It\^o's formula without growth conditions}\label{sec:unrestricted}

The $L^2$ formulas of \cref{sec:ito} need the integrand $f_x(\cdot,B)$ to lie in $\HT$, which is what the
derivative bounds and growth bounds secure. Without them the stochastic term is no longer an $L^2$ integral,
and the classical statement is that the compensated process is a continuous local martingale. That is the
statement proved here.

\begin{theorem}[\lean{ito_formula_unrestricted}]\label{thm:unrestricted}
Let $B$ have measurable coordinates and continuous paths, and let $f\colon\R\times\R\to\R$ have everywhere
partial derivatives $f_t,f_x,f_{xx},f_{tt},f_{tx},f_{xxx}$ which, together with $f$, are jointly continuous.
No growth or boundedness condition is imposed. Then there is a process $M$ with every path continuous, which
is a local martingale for the null-augmented filtration $(\F^\mu_t)$, such that for every $t\ge0$, almost
surely,
\[
  f(t,B_t)-f(0,B_0)=M_t+\int_{(0,t]}\bigl(f_t+\tfrac12f_{xx}\bigr)(s,B_s)\,\dd s .
\]
\end{theorem}

\begin{leancode}[caption={It\^o's formula without growth conditions (\texttt{Foundations/ItoFormulaUnrestrictedLocMart}); five of the six derivative hypotheses and five of the seven continuity hypotheses are abbreviated.},label=lst:unrestricted]
theorem ito_formula_unrestricted (hB : IsPreBrownianReal B μ)
    (hBmeas : ∀ t, Measurable (B t)) (hBcont : ∀ ω, Continuous fun s : ℝ≥0 ↦ B s ω)
    {f f_t f_x f_xx f_tt f_tx f_xxx : ℝ → ℝ → ℝ}
    (hf_t : ∀ t x, HasDerivAt (fun s ↦ f s x) (f_t t x) t)
    -- ... hf_tt, hf_tx, hf_x, hf_xx, hf_xxx: the other five partial derivatives
    (hf_cont : Continuous fun p : ℝ × ℝ ↦ f p.1 p.2)
    -- ... joint continuity of f_t, f_x, f_xx, f_tt, f_tx
    (hf_xxx_cont : Continuous fun p : ℝ × ℝ ↦ f_xxx p.1 p.2) :
    ∃ M : ℝ≥0 → Ω → ℝ,
      (∀ ω, Continuous fun t ↦ M t ω) ∧
      IsLocalMartingale M (augFiltration (μ := μ) hBmeas) μ ∧
      (∀ t : ℝ≥0, (fun ω ↦ f (t : ℝ) (B t ω) - f 0 (B 0 ω)) =ᵐ[μ]
        (fun ω ↦ M t ω + ∫ s in Set.Ioc 0 t,
          (f_t (s : ℝ) (B s ω) + (1 / 2) * f_xx (s : ℝ) (B s ω)) ∂timeMeasure))
\end{leancode}

\subsection{Exit times as localizers}
The localizing sequence is the exit time of the closed exterior $\{\abs{x}\ge N\}$,
$\tau_N=\inf\{t:\abs{B_t}\ge N\}$ (\lean{exitTime}). Because the hit set is closed and the paths continuous,
$\{\tau_N\le i\}$ is the event $\{\exists s\in[0,i]:\abs{B_s}\ge N\}$, which is measurable for $\F_i$ through a
countable rational description. So $\tau_N$ is a stopping time for the raw natural filtration, with no
right-continuity required; the open exterior $\{\abs{x}>N\}$ would characterise only $\{\tau_N<i\}$, an event of
$\F_{i+}$. The package's own localizing sequence of hitting times needs a complete and right-continuous
filtration, which the natural filtration is not; the exit time does not. For continuous paths the exit times
escape to infinity, since each path is bounded on compacts.

\subsection{Double truncation and agreement on stochastic intervals}
Localizing in space is not enough: a general $f$ has $t$-derivatives unbounded over $t\in\R$. The proof therefore
truncates both arguments, $f^N(t,x)=f(\mathrm{cut}_N(t),\mathrm{cut}_N(x))$, with a smooth cutoff onto a compact
square on which the continuous partials of $f$ are bounded. The chain rule gives the six partials of $f^N$ in
closed form, each globally bounded, so the process form of the bounded formula applies and yields a truncated
It\^o identity with a continuous \emph{true}-martingale modification $M^N$ on the half-line. On the stochastic
interval $\{t\le\sigma_N\}$ with $\sigma_N=\min(\tau_N,N)$ both cuts are inert ($t\le N$ and $\abs{B_s}\le N$ for
$s\le t$), so $f^N=f$ there and $M^N_t=M_t$ almost surely on $\{t\le\sigma_N\}$. Capping the localizer at $N$ is
what keeps the time cut inert exactly where the agreement is used.

\subsection{From per-time agreement to a local martingale}
Agreement at each fixed $t$ is lifted to agreement of whole paths on the stochastic interval, almost surely,
by continuity and a countable dense set of times (\lean{indistinguishable_on_stochInterval}). The
compensated process $M$ is adapted, because its drift primitive $\int_0^t(f_t+\frac12f_{xx})(s,B_s)\,\dd s$
is $\F_t$-measurable: after clamping the integrand's time to $[0,t]$ the integrand is a Carath\'eodory
function (continuous in $s$, $\F_t$-measurable in $\omega$), hence jointly measurable, and its integral is
measurable (\lean{StronglyMeasurable.integral_prod_right}). With $M$ adapted, the stopped indicator processes
are martingales, and they assemble into the package's \lean{IsLocalMartingale} with localizing sequence
$\sigma_N$.

\subsection{What the theorem says, and what it does not}
\Cref{thm:unrestricted} removes the growth conditions of \cref{sec:ito} entirely, but it says less than the
classical formula in one respect: $M$ is not identified as a stochastic integral $\int_0^tf_x(s,B_s)\,\dd B_s$,
because the library's integral is defined for square-integrable integrands and $f_x(\cdot,B)$ need not be one.
Where it is, the identification is \cref{thm:ito-growth}. The theorem covers every $C^3$ function, including
those that grow faster than any exponential: applied to $f(t,x)=e^{x^2}$, for instance, it gives a continuous
local martingale with no integrability of $f(t,B_t)$ assumed. The smoothness hypothesis is the one the bounded
engine needs; the $C^{1,2}$ statement would require an additional approximation of $f$ and is not formalized.

%======================================================================
\section{It\^o processes}\label{sec:itoprocess}

Every model in continuous-time finance is an It\^o process $X=x_0+\int b\,\dd s+\int\sigma\,\dd B$, so the
formula must reach functions of $X$, not only of $(t,B_t)$. Three results lead there.

\paragraph{Quadratic variation of an It\^o process}
For $X=X_0+A+\sigma\bulletB B$ with a drift path $A$ Lipschitz in time uniformly in $\omega$, the squared-increment
sums along the uniform partition converge to $\int_0^T\sigma_s^2\,\dd s$: in $L^2$ for constant $\sigma$
(\file{ItoProcessQV}), and in $L^1$ for $\sigma$ bounded, adapted and continuous in time, with the weighted
form $\sum_kw(t_k)(\Delta X_k)^2\to\int_0^Tw\sigma^2\,\dd s$ for bounded adapted weights with almost surely
continuous paths. The drift contributes nothing. The proof splits $(\Delta X_k)^2$ around the frozen increment
$\sigma(t_k)\Delta B_k$: the freezing defects $D_k=\Delta M_k-\sigma(t_k)\Delta B_k$ of $M=\sigma\bulletB B$ are
orthogonal It\^o integrals whose sum is the error of the frozen Riemann sum, so $\sum_k\norm{D_k}^2\to0$
(\file{AdaptedStochasticIntegralFreezing}), and what remains is the weighted quadratic variation of $B$. The
$L^1$ mode is a limitation of the method: an $L^2$ bound on the freezing part would need fourth moments of
$\sigma\bulletB B$.

\paragraph{Riemann--Stieltjes sums against an It\^o integral}
For $M=\varphi\bulletB B$ and a bounded predictable weight $w$ whose paths are almost surely continuous on $[0,T]$,
the sums $\sum_kw(t_k)(M_{t_{k+1}}-M_{t_k})$ converge in mean square to the integral of $w$ against $M$
(\cref{sec:against}), which is the It\^o integral against $B$ of the class of $\varphi w$
(\file{AdaptedRiemannStieltjes}). The route goes through $L^2(\langle M\rangle)$: each sum is the integral
against $M$ of a left-endpoint step process, the step processes converge to $w$ by dominated convergence, and
the isometry against $M$ carries the convergence over. No path regularity of $\varphi$ is needed.

\paragraph{Constant coefficients}
For $X_t=X_0+bt+\sigma B_t$ and $f$ whose first three derivatives grow at most exponentially, the formula
$f(X_T)-f(X_0)=\ito_T(\sigma f'(X))+\int_0^T(f'(X_s)b+\frac12f''(X_s)\sigma^2)\,\dd s$ holds with the integrand
named (\lean{ito_formula_itoProcess}); this is the form the geometric Brownian motion needs.

\begin{theorem}[\lean{ito_formula_adapted}]\label{thm:adapted}
Let $B$ have every path continuous, let $\sigma$ and $b$ be bounded, adapted to the natural filtration of $B$,
with every path continuous, and let $X=x_0+\int_0^\cdot b\,\dd s+\int_0^\cdot\sigma\,\dd B$. For $f$ twice
differentiable with $f'$, $f''$ bounded and $f''$ continuous, almost surely
\[
  f(X_T)-f(x_0)=\ito_T\bigl(f'(X)\,\sigma\bigr)+\int_0^T\Bigl(f'(X_s)\,b_s+\tfrac12f''(X_s)\,\sigma_s^2\Bigr)\dd s ,
\]
the stochastic integral being the It\^o integral at horizon $T$ of the class of the bounded predictable process
$f'(X)\sigma$.
\end{theorem}

The process $X$ is defined (\lean{adaptedItoProcess}) as $x_0$ plus the pathwise drift integral plus the
continuous modification of the It\^o integral process of $\sigma$. The proof applies the discrete It\^o formula
on the uniform partition, which splits $f(X_T)-f(x_0)$ into drift sums, Riemann--Stieltjes sums against
$M=\sigma\bulletB B$, weighted quadratic variation and a remainder. The three random limits hold in measure; a
common subsequence along which all three converge almost surely is extracted, and on such a path, where $X$ is
also continuous, a deterministic lemma (\lean{ito_formula_of_tendsto}) passes the discrete identity to the
limit. Writing the stochastic term as an integral against $M$ and then, by the chain rule of
\cref{sec:against}, against $B$ gives the displayed form. Because the limit is taken almost surely along a
subsequence rather than in $L^2$ at a rate, the Taylor remainder is controlled by the modulus of continuity of $f''$
times the quadratic variation, and no third derivative is needed: this is the library's first It\^o formula for
$C^2$ functions, as well as its first for a process that is not a function of $(t,B_t)$. Its scope is recorded in
\cref{sec:scope}.

%======================================================================
\section{The integral as an operator: its range and its uniqueness}\label{sec:operator}

Two theorems characterise $\ito_T$ as an operator. Its range is as large as it can be, and it is the only
reasonable integral.

\subsection{Martingale representation}\label{sec:representation}

\begin{theorem}[\lean{exists_itoIntegral_representation}]\label{thm:mrt}
Let $B$ have every path continuous. For every $F\in L^2(\mu)$ that is $\F^B_T$-measurable there is a unique
$\varphi\in\HT$ with
\[
  F=\E[F]+\ito_T(\varphi)\qquad\mu\text{-a.s.}
\]
Equivalently, the range of $\ito_T$ together with the constants is the whole of $L^2(\F^B_T)$
(\lean{itoIntegralCLM_T_surjective_onto_centered}), and $\ito_T$ is an isometric isomorphism of $\HT$ onto the
centred part of $L^2(\F^B_T)$ (\lean{itoIsometryEquiv}).
\end{theorem}

\begin{leancode}[caption={Martingale representation, terminal form (\texttt{Foundations/MartingaleRepresentation}).},label=lst:mrt]
theorem exists_itoIntegral_representation (hBmeas : ∀ t, Measurable (B t))
    (hBcont : ∀ ω, Continuous fun s : ℝ≥0 ↦ B s ω) (T : ℝ≥0)
    (F : Lp ℝ 2 μ)
    (hFmeas : AEStronglyMeasurable[ItoIntegralL2.natFiltration hBmeas T] (⇑F) μ) :
    ∃! φ : Lp ℝ 2 (trimMeasure_T (μ := μ) T hBmeas),
      (⇑F) =ᵐ[μ] fun ω ↦ (∫ x, F x ∂μ) + itoIntegralCLM_T hB T hBmeas φ ω
\end{leancode}

The process form follows from \cref{thm:key}: every square-integrable martingale for the natural filtration is
$M_t=M_0+(\varphi\bulletB B)_t$ for $t\le T$ (\lean{martingale_representation}), because $(\varphi\bulletB B)_t$ is the
$\F_t$-conditional expectation of $\ito_T(\varphi)$.

The proof uses no Malliavin calculus, no Clark--Ocone formula and no It\^o formula for an adapted integrand. The
range of $\ito_T$ is closed, being the image of an isometry, so it suffices to show that a centred $\xi\in L^2(\F^B_T)$
orthogonal to the range is zero. Two facts do this. First, for every step function $h$ the Dol\'eans exponential
$\mathcal E(h)=\exp(\int h\,\dd B-\frac12\int h^2\,\dd s)$ satisfies $\mathcal E(h)-1\in\operatorname{range}\ito_T$
(\lean{stepDoleans_sub_one_mem_range}), starting from the constant step, which is the discounted geometric Brownian
motion of \cref{sec:growth}, and adding one step at a time with the locality of the integral. Second, an element
orthogonal to every $\mathcal E(h)$ vanishes (\lean{eq_zero_of_orthogonal_stepDoleans}). Abel summation turns the step
exponentials into the exponentials $\exp(\sum_i\lambda_iB_{t_i})$ of finitely many coordinates; on each dyadic
cylinder $\sigma$-algebra the conditional expectation of an orthogonal element then has a vanishing Laplace transform,
which forces it to zero by analytic continuation of the complex moment-generating function and uniqueness of
characteristic functions; and by path continuity the dyadic cylinder $\sigma$-algebras increase to $\F^B_T$
(\lean{iSup_cylinderFiltration_eq_natFiltration}), so L\'evy's upward theorem sends the conditional expectations to
$\xi$ itself. Centring,
$\E[\ito_T(\varphi)]=0$, is proved a floor down (\lean{integral_itoIntegralCLM_T}).

In the parent library this theorem is market completeness: every square-integrable $\F^B_T$-measurable claim is
attained from initial wealth $\E[H]$ by a unique strategy (\lean{exists_replicating_strategy}). What is not
proved is the explicit form of $\varphi$, which is the Clark--Ocone formula and needs the Malliavin derivative.

\subsection{Integration against an It\^o integral}\label{sec:against}
For $\varphi\in\HT$ let $M=\varphi\bulletB B$. The bracket of $M$ is the measure $\langle M\rangle=\varphi^2\cdot(\lambda_T\otimes\mu)$
on the predictable $\sigma$-algebra (\lean{bracketMeasure}), and the integral against $M$ is defined on
$L^2(\langle M\rangle)$ by composing $\ito_T$ with the multiplication isometry $\psi\mapsto\psi\varphi$ from
$L^2(\langle M\rangle)$ to $\HT$:
\[
  \int\psi\,\dd M=\ito_T(\psi\varphi),\qquad\Bigl\lVert\int\psi\,\dd M\Bigr\rVert_{L^2(\mu)}=\norm{\psi}_{L^2(\langle M\rangle)}
\]
(\lean{itoIntegralAgainstCLM}, \lean{norm_itoIntegralAgainstCLM}). The chain rule is thus how the object is built; Mathlib has the multiplication isometry only for $p=1$, so the
$p=2$ case is proved here (\file{LpMulIsometry}). What identifies the result as \emph{the} integral against $M$ is the
Riemann--Stieltjes value on elementary integrands, $\int Z\ind_{(a,b]}\,\dd M=Z(M_b-M_a)$ for bounded
$\F_a$-measurable $Z$ (\lean{itoIntegralAgainst_elementary}), together with uniqueness among continuous linear maps
with those values (\lean{itoIntegralAgainst_unique}).

The bracket earns its name. The second moment of an increment is its mass,
$\E[(M_b-M_a)^2]=\langle M\rangle((a,b]\times\Omega)$ (\lean{norm_sq_increment_eq_bracket}); conditionally,
$\E[(M_b-M_a)^2\mid\F_a]=\E[\int_a^b\varphi_u^2\,\dd u\mid\F_a]$ against the pathwise increment of the bracket
(\lean{condExp_band_second_moment}); and that increment is adapted (\lean{bracketProcess_adapted}), so
$\E[M_b^2-\langle M\rangle_b\mid\F_a]=M_a^2-\langle M\rangle_a$ (\lean{condExp_sq_sub_bracket}): the bracket is the
compensator of $M^2$. Not formalized are any
pathwise quadratic variation of $M$, a bundled martingale structure for $M^2-\langle M\rangle$, and the Doob--Meyer
decomposition of a general submartingale.

\subsection{Agreement with the upstream characterisation}\label{sec:characterisation}
\texttt{BrownianMotion} characterises a stochastic integral axiomatically~\cite{BrownianMotionRepo}. A map $I$ from
processes to random variables, defined on a domain $S$, is a \emph{Riemann--Stieltjes extension} if it returns
$X(Y_j-Y_i)$ on every elementary process $\ind_{(i,j]}X$ with simple $\F_i$-measurable $X$, is linear, respects
indistinguishability, and passes to the limit under dominated convergence; it is a \emph{stochastic integral} if,
in addition, it agrees almost surely with every Riemann--Stieltjes extension on the intersection of their domains.

\begin{theorem}[\lean{isStochasticIntegral}]\label{thm:char}
For every $\varphi\in\HT$, the integral against $M=\varphi\bulletB B$, with integrator $M$ stopped at $T$, the natural
filtration, and as domain the processes indistinguishable from a predictable process square-integrable against
$\langle M\rangle$, is a stochastic integral in the sense of \texttt{BrownianMotion}.
\end{theorem}

\begin{leancode}[caption={The characterisation (\texttt{Foundations/StochasticIntegralCharacterisation}).},label=lst:char]
theorem isStochasticIntegral (hB : IsPreBrownianReal B μ) :
    IsStochasticIntegral μ (integrator T hBmeas φ hB) (natFiltration hBmeas)
      (sIntegral T hBmeas φ hB) (domain T hBmeas φ)
\end{leancode}

Taking $\varphi=1$ makes $M$ the Brownian motion stopped at $T$, so the theorem says in particular that $\ito_T$ is a
stochastic integral against $B$ in the upstream sense. Two features of the proof are worth recording. The domain is forced
rather than chosen: the predicate asks for closure under indistinguishability, and a process indistinguishable
from a predictable one need not be predictable, or even jointly measurable, so the domain is the processes with a
predictable square-integrable version, and the integral reads the class of any version. And the uniqueness clause is
the substantive part. An extension in the upstream sense is not assumed to be $L^2$-continuous, so uniqueness among
continuous linear maps does not reach it; the proof is a monotone-class argument from dominated convergence alone,
through a Dynkin argument over the predictable rectangles, then simple functions, bounded processes, and finally the
whole common domain by truncation. Dominated convergence in $L^p$, which Mathlib at the pin has only for $p=1$, is
proved on the way for every $p<\infty$ and every measure (\lean{tendsto_eLpNorm_sub_of_dominated_convergence}).

%======================================================================
\section{What the calculus carries}\label{sec:consumers}

The integral and the formula are consumed downstream. Two consumers belong to stochastic analysis proper.

\paragraph{Girsanov's theorem}
For a bounded predictable drift $\theta$, $\abs{\theta}\le C$, let
$\frac{\dd Q}{\dd\mu}=\exp\bigl(-\ito_T(\theta)-\frac12\int_0^T\theta_s^2\,\dd s\bigr)$. Under $Q$ the process
$B^\theta_t=B_t+\int_0^t\theta_s\,\dd s$ starts at $0$, has increments $B^\theta_t-B^\theta_s\sim N(0,t-s)$ for
$s\le t\le T$, and has jointly independent increments on $[0,T]$
(\lean{Btheta_isQBrownianMotion_predictable_of_bdd}); these fix every finite-dimensional law of $B^\theta$ on
$[0,T]$ to that of a Brownian motion. Path continuity of $B^\theta$ and independence of its increments from the
past are not stated. The theorem is the last rung of a ladder (constant, simple adapted, continuous adapted, bounded
predictable $\theta$) whose rungs share one lemma: any process supplying its own exponential martingale under $Q$
has the Brownian finite-dimensional laws (\lean{isQBrownianMotion_of_expMartingale}), by characteristic functions,
with no adapted-integrand It\^o formula. The bounded-predictable statement is made for an adapted modification of
$B^\theta$, which has the same finite-dimensional laws. The general case under Novikov's condition is not
formalized.

\paragraph{Stochastic differential equations}
For Lipschitz $b$ and $\sigma$, with constants $L_b$ and $L_\sigma$, the Picard map
$X\mapsto\eta+\int_0^\cdot b(X)\,\dd s+\int_0^\cdot\sigma(X)\,\dd B$ is a contraction of $\HT$ with constant
$TL_b+\sqrt T\,L_\sigma$, so when that constant is below $1$ it has a unique fixed point
(\lean{picardMap_exists_unique_fixedPoint}). The diffusion term is the It\^o integral of this paper; the fixed point
is realized as a sample-path process, $X_t=\eta+\int_0^tb(X_s)\,\dd s+(\sigma(X)\bulletB B)_t$ with the drift identified
as a pathwise Lebesgue integral (\lean{sde_pathwise_decomposition},
\lean{driftContinuousMod_eq_setIntegral}). Pathwise uniqueness of $L^2$ strong solutions
follows from an energy estimate and Gr\"onwall's inequality (\lean{sde_pathwise_uniqueness}); in that statement the
diffusion term is an abstract operator whose only assumed property is the It\^o isometry bound. Existence on
arbitrary horizons, which needs a weighted norm, is not formalized.

%======================================================================
\section{Scope: what is, and is not, an It\^o calculus here}\label{sec:scope}

The title claims ``a'' calculus, and the indefinite article is doing deliberate work. What earns the noun is that the
canonical objects and theorems are present and \emph{compose}: the stochastic term of It\^o's formula is the isometric
integral of \cref{sec:integral}, whose energy law is the isometry of \cref{sec:process}, whose conditional projection is
the martingale of \cref{thm:key}, whose continuous representative is the local martingale of \cref{sec:pathwise}, whose
range is characterised by \cref{thm:mrt}, and which is the upstream package's stochastic integral by \cref{thm:char}. One
object, a single chain of theorems, no re-derivation gluing them. What the indefinite article concedes is equally
concrete.

\paragraph{The integrator}
The integrator is Brownian motion, or an It\^o integral against it. Integration against a general continuous local
martingale or semimartingale is not formalized; \cref{thm:char} places the present integral inside the upstream
package's axiomatic characterisation of a stochastic integral.

\paragraph{Square-integrable integrands}
Integrands live in $L^2$: in $\HT$ on $[0,T]$, or in its unbounded-horizon analogue. The localized class
$\{\varphi:\int_0^T\varphi_s^2\,\dd s<\infty\text{ a.s.}\}$ is not covered, which is why the local martingale of
\cref{thm:unrestricted} is not identified as a stochastic integral.

\paragraph{The scope of each It\^o formula}
In $L^2$, the formula for $f(t,B_t)$ needs the six partial derivatives bounded or of exponential growth in $x$. In
local-martingale form it needs them continuous and nothing else. For It\^o processes with adapted coefficients it is
proved at a fixed time $T$, for bounded coefficients with every path continuous, and for $f$ with bounded $f'$ and
$f''$; so it does not yet cover $e^x$ or $x^2$ of such a process, time-dependent $f$, a random initial value, or
coefficients of a stochastic differential equation. The multi-dimensional formula is recorded in the library only as
a reduced core, whose hypotheses include its conclusion.

\paragraph{Completeness, not the usual conditions}
The pathwise local martingales live on the null-augmented filtration: completeness, which the continuous modification
needs, is added, but right-continuity, the other half of the usual conditions, is neither assumed nor used.

\paragraph{Quadratic variation in mean, not pathwise}
Quadratic variation is established in $L^1$ and $L^2$ along uniform or refining partitions, for Brownian motion and
for It\^o processes, and the bracket of $\varphi\bulletB B$ is identified in conditional mean. No pathwise (almost sure)
quadratic variation is formalized, and L\'evy's characterisation of Brownian motion is recorded only as a reduced
core.

None of these boundaries qualifies a statement made in \cref{sec:integral,sec:process,sec:ito,sec:pathwise,sec:unrestricted,sec:itoprocess,sec:operator,sec:consumers};
they bound the \emph{reach} of the development, not the \emph{validity} of any theorem in it.

%======================================================================
\section{Related work}\label{sec:related}

Keskin's Isabelle/HOL martingale library~\cite{Keskin2023} is the closest prior art in the probabilistic direction;
it develops sub- and supermartingales in discrete and continuous time and stops below stochastic integration. Degenne,
Ledvinka, Marion and Pfaffelhuber~\cite{DegenneEtAl2025,BrownianMotionRepo} construct Brownian motion in Lean, through
the Kolmogorov extension theorem, Gaussian measures and a Kolmogorov--Chentsov theorem; their package is the foundation
of this work, and its stochastic-integration programme (an axiomatic characterisation of the integral, local
martingales, quadratic variation of locally square-integrable martingales, and the Doob--Meyer decomposition) is
complementary to it: a general theory under construction, against which \cref{thm:char} certifies the concrete
Brownian integral built here. Nagy's Lean derivation from It\^o's lemma to Black--Scholes~\cite{Nagy2026} formalizes the
structural drift computation of It\^o's lemma and, on its own account, leaves the continuous It\^o integral
structurally verified rather than constructed. Avigad, H\"olzl and Serafin's central limit theorem~\cite{AvigadHolzlSerafin2017}
remains the historical anchor for deep measure-theoretic probability in proof assistants, and Echenim, Guiol and
Peltier formalized discrete-time replication pricing in Isabelle/HOL~\cite{EchenimGuiolPeltier2020}. StochBench~\cite{StochBench2026}
is a recent Lean benchmark of graduate stochastic-process problems which, lacking infrastructure such as the It\^o
integral, states many of its targets with the required properties as hypotheses. A search in October 2026 of the
Isabelle Archive of Formal Proofs and the Lean and Coq/Rocq ecosystems found no other completed machine-checked
construction of the It\^o integral, proof of It\^o's formula, or martingale representation theorem.

%======================================================================
\section{Conclusion}\label{sec:conclusion}

The development formalizes the It\^o calculus of Brownian motion: on $[0,T]$ the integral as an isometry, the integral as
a continuous $L^2$ martingale with its full second-moment structure, its range (martingale representation) and its
uniqueness (agreement with the upstream axiomatic integral); It\^o's formula in $L^2$ under bounded or exponentially
growing derivatives, as a process, and for It\^o processes with adapted coefficients; and the pathwise object, a
continuous local martingale on the half-line, to which It\^o's formula extends without any growth condition. To our
knowledge it is the first machine-checked It\^o calculus in any proof assistant. It is also the layer the parent
library~\cite{CoelhoLibrary} stands on, where the integral and the formula are consumed to construct, rather than assume,
objects of continuous-time pricing: the dynamic change of measure, the replicating hedge, and the discounted price
as an It\^o integral. The boundaries of \cref{sec:scope} (general semimartingale integrators, the
localized integrand class, pathwise quadratic variation, and the multi-dimensional formula) mark where the formalized
theory currently ends.

\paragraph{Provenance}
The Lean development was written by the author with an AI coding assistant (Claude Code), as recorded in the
repository's provenance file; every proof is checked by the Lean kernel and passes the gates of \cref{sec:setting}
unchanged. This manuscript was revised with the same assistant, and its statements were checked against the Lean
sources of release \texttt{v1.5.0}.

%======================================================================
\appendix
\section{Displayed results and their Lean names}\label{app:names}

\begin{center}
\scriptsize
\begin{tabularx}{\linewidth}{@{}>{\raggedright\arraybackslash}p{0.21\linewidth}>{\raggedright\arraybackslash}p{0.33\linewidth}>{\raggedright\arraybackslash}X@{}}
\toprule
Result & Lean name & Module / benchmark entry \\
\midrule
Rectangle $\pi$-system & \lean{isPiSystem_predictableRect} & \file{ItoIntegralCLM} \\
Generation & \lean{generateFrom_predictableRect} & \file{ItoIntegralCLM} \\
Density & \lean{simpleAssembly_T_denseRange} & \file{ItoIntegralCLM} \\
The It\^o integral, isometry & \lean{itoIntegralCLM_T}, \lean{itoIntegralCLM_T_norm} & \file{ItoIntegralCLM} \\
Unbounded horizon & \lean{itoIntegralL2_norm} & \file{ItoIntegralL2Dense} \\
Keystone $\ito_T(g_B)$ & \lean{itoIntegralCLM_T_brownian} & \file{ItoIntegralBrownian} \\
Conditional increments & \lean{condExp_adapted_mul_increment} & \file{ItoIntegralProcessMartingale} \\
Key identity (\cref{thm:key}) & \lean{itoProcessCLM_eq_condExpL2} & \file{ItoIntegralProcessGeneral} \\
Martingale property & \lean{itoIntegralProcessGen_isMartingale} & \file{ItoIntegralProcessGeneral} \\
Time-indexed isometry & \lean{itoProcessCLM_norm_sq} & \file{ItoIntegralProcessIsometry} / \bench{sc-ito-general-time-isometry} \\
QV in $L^2$, in $L^1$ & \lean{tendsto_qv}, \lean{qv_equals_t} & \file{QuadraticVariationL2}, \file{BrownianQuadraticVariation} / \bench{sc-thm-6.1.1} \\
Weighted QV & \lean{tendsto_weighted_qv_process} & \file{WeightedQuadraticVariation} \\
It\^o, bounded (\cref{thm:ito-bdd}) & \lean{ito_formula_L2_bddDeriv} & \file{ItoFormulaCLM} / \bench{sc-thm-7.1.1} \\
It\^o, time-dependent (\cref{thm:ito-td}) & \lean{ito_formula_td_L2_bddDeriv} & \file{ItoFormulaTD} / \bench{sc-thm-7.1.2} \\
It\^o as a process & \lean{ito_formula_td_process} & \file{ItoFormulaProcess} / \bench{sc-ito-formula-td-process} \\
It\^o, exponential growth (\cref{thm:ito-growth}) & \lean{ito_formula_td_localized} & \file{ItoFormulaLocalized} / \bench{sc-ito-formula-localized} \\
GBM; discounted GBM & \lean{ito_formula_gbm}, \lean{discountedGBM_eq_itoIntegral} & \file{ItoFormulaGBM} / \bench{sc-ito-formula-gbm}, \bench{sc-discounted-gbm-ito} \\
Expectation It\^o & \lean{expectation_ito} & \file{FeynmanKacHeatEquation} \\
Continuous modification & \lean{exists_continuous_modification_itoProcess} & \file{ItoIntegralProcessContinuousModification} \\
Local martingale on $[0,T]$ & \lean{exists_continuous_localMartingale_modification} & \file{ItoIntegralProcessLocalMartingaleGeneral} \\
Local martingale on $\R_{\ge0}$ & \lean{exists_continuous_localMartingale_modification_infinite} & \file{ItoIntegralProcessLocalMartingaleInfinite} / \bench{sc-ito-infinite-local-martingale} \\
It\^o, no growth (\cref{thm:unrestricted}) & \lean{ito_formula_unrestricted} & \file{ItoFormulaUnrestrictedLocMart} / \bench{sc-ito-formula-unrestricted-islocalmartingale} \\
Exit times & \lean{exitTime} & \file{ExitTime} \\
It\^o, constant coefficients & \lean{ito_formula_itoProcess} & \file{ItoFormulaItoProcess} / \bench{sc-ito-formula-ito-process} \\
It\^o, adapted coefficients (\cref{thm:adapted}) & \lean{ito_formula_adapted} & \file{ItoFormulaAdapted} / \bench{sc-ito-formula-adapted} \\
Representation (\cref{thm:mrt}) & \lean{exists_itoIntegral_representation}, \lean{itoIntegralCLM_T_surjective_onto_centered}, \lean{itoIsometryEquiv} & \file{MartingaleRepresentation} \\
Representation, process form & \lean{martingale_representation} & \file{MartingaleRepresentation} / \bench{gir-thm-9.3.4} \\
Integral against $\varphi\bulletB B$ & \lean{itoIntegralAgainstCLM}, \lean{itoIntegralAgainst_elementary} & \file{ItoIntegralAgainstMartingale} / \bench{sc-ito-integral-band} \\
Bracket, conditional & \lean{condExp_band_second_moment} & \file{PointwiseBracket} / \bench{sc-bracket-conditional} \\
Characterisation (\cref{thm:char}) & \lean{isStochasticIntegral} & \file{StochasticIntegralCharacterisation} / \bench{sc-ito-is-stochastic-integral} \\
Girsanov, bounded predictable & \lean{Btheta_isQBrownianMotion_predictable_of_bdd} & \file{GirsanovPredictableTheta} / \bench{gir-thm-9.1.8-predictable} \\
SDE existence and uniqueness & \lean{picardMap_exists_unique_fixedPoint} & \file{SDEExistence} / \bench{sde-picard-existence-uniqueness} \\
\bottomrule
\end{tabularx}
\end{center}
All modules are under \file{MathFin/Foundations/}. The library contains no \texttt{sorry}. Every theorem in the table is pinned
directly by one of the two axiom audits to the three standard axioms, except the supporting lemma \lean{isPiSystem_predictableRect},
which is used in the proofs of pinned theorems; definitions are not pinned separately.

%======================================================================
\begin{thebibliography}{99}
\small
\bibitem{AvigadHolzlSerafin2017} J.~Avigad, J.~H\"olzl, L.~Serafin. A formally verified proof of the central limit theorem.
  \emph{Journal of Automated Reasoning}, 59(4):389--423, 2017.
\bibitem{BrownianMotionRepo} R.~Degenne et al. \emph{brownian-motion}: a Lean~4 development of Brownian motion and stochastic
  integrals. \url{https://github.com/RemyDegenne/brownian-motion}, 2024--2026. Consumed at commit \texttt{0d5b6eb}.
\bibitem{CoelhoFTAP} R.~Coelho. The fundamental theorem of asset pricing, formalized in Lean~4. arXiv:2606.28990, 2026.
\bibitem{CoelhoLibrary} R.~Coelho. A formally verified library of mathematical finance in Lean~4. arXiv:2606.01356, 2026.
\bibitem{DegenneEtAl2025} R.~Degenne, D.~Ledvinka, E.~Marion, P.~Pfaffelhuber. Formalization of Brownian motion in Lean.
  arXiv:2511.20118, 2025.
\bibitem{deMouraUllrich2021} L.~de~Moura, S.~Ullrich. The Lean~4 theorem prover and programming language. In \emph{CADE~28},
  LNCS 12699, pp.~625--635, 2021.
\bibitem{EchenimGuiolPeltier2020} M.~Echenim, H.~Guiol, N.~Peltier. Formalizing the Cox--Ross--Rubinstein pricing of European
  derivatives in Isabelle/HOL. \emph{Journal of Automated Reasoning}, 64:737--765, 2020.
\bibitem{Ito1944} K.~It\^o. Stochastic integral. \emph{Proceedings of the Imperial Academy, Tokyo}, 20(8):519--524, 1944.
\bibitem{KaratzasShreve1991} I.~Karatzas, S.~E.~Shreve. \emph{Brownian Motion and Stochastic Calculus}, 2nd ed. Springer, 1991.
\bibitem{Keskin2023} A.~Keskin. A formalization of martingales in Isabelle/HOL. arXiv:2311.06188, 2023.
\bibitem{Mathlib2020} The mathlib Community. The Lean mathematical library. In \emph{CPP 2020}, pp.~367--381, 2020.
\bibitem{Nagy2026} T.~Nagy. From It\^o to Black--Scholes: a machine-verified derivation in Lean~4. SSRN Working Paper 6336503,
  2026.
\bibitem{StochBench2026} I.~Davidovich, D.~Ganguly, V.~Singh, V.~Chaudhary. StochBench: a domain-specific benchmark for
  stochastic processes in Lean. arXiv:2609.09264, 2026.
\end{thebibliography}

\end{document}
